\documentclass[10pt,a4paper]{amsart}
\usepackage{amsmath,amssymb,amsthm,graphicx}
% Retained minimal preamble for the Stacks Project chapter more-morphisms.tex (More on Morphisms).
% Verbatim excerpt of the chapter preamble: Xy-pic setup (native preamble lines 14-18) and the native
% theorem declarations with the mathematical macros (native preamble lines 154-217). Copied mechanically;
% nothing added, removed or redefined. The wrapper class and the amsmath/amssymb/amsthm/graphicx base
% packages are supplied by the assembler, so the class/externaldocument part of the native preamble is omitted.
% For commutative diagrams we use Xy-pic
\usepackage[all]{xy}

% We use 2cell for 2-commutative diagrams.
\xyoption{2cell}
% Theorem environments.
%
\theoremstyle{plain}
\newtheorem{theorem}[subsection]{Theorem}
\newtheorem{proposition}[subsection]{Proposition}
\newtheorem{lemma}[subsection]{Lemma}

\theoremstyle{definition}
\newtheorem{definition}[subsection]{Definition}
\newtheorem{example}[subsection]{Example}
\newtheorem{exercise}[subsection]{Exercise}
\newtheorem{situation}[subsection]{Situation}

\theoremstyle{remark}
\newtheorem{remark}[subsection]{Remark}
\newtheorem{remarks}[subsection]{Remarks}

\numberwithin{equation}{subsection}

% Macros
%
\def\lim{\mathop{\mathrm{lim}}\nolimits}
\def\colim{\mathop{\mathrm{colim}}\nolimits}
\def\Spec{\mathop{\mathrm{Spec}}}
\def\Hom{\mathop{\mathrm{Hom}}\nolimits}
\def\Ext{\mathop{\mathrm{Ext}}\nolimits}
\def\SheafHom{\mathop{\mathcal{H}\!\mathit{om}}\nolimits}
\def\SheafExt{\mathop{\mathcal{E}\!\mathit{xt}}\nolimits}
\def\Sch{\mathit{Sch}}
\def\Mor{\mathop{\mathrm{Mor}}\nolimits}
\def\Ob{\mathop{\mathrm{Ob}}\nolimits}
\def\Sh{\mathop{\mathit{Sh}}\nolimits}
\def\NL{\mathop{N\!L}\nolimits}
\def\CH{\mathop{\mathrm{CH}}\nolimits}
\def\proetale{{pro\text{-}\acute{e}tale}}
\def\etale{{\acute{e}tale}}
\def\QCoh{\mathit{QCoh}}
\def\Ker{\mathop{\mathrm{Ker}}}
\def\Im{\mathop{\mathrm{Im}}}
\def\Coker{\mathop{\mathrm{Coker}}}
\def\Coim{\mathop{\mathrm{Coim}}}

% Boxtimes
%
\DeclareMathSymbol{\boxtimes}{\mathbin}{AMSa}{"02}

%
% Macros for moduli stacks/spaces
%
\def\QCohstack{\mathcal{QC}\!\mathit{oh}}
\def\Cohstack{\mathcal{C}\!\mathit{oh}}
\def\Spacesstack{\mathcal{S}\!\mathit{paces}}
\def\Quotfunctor{\mathrm{Quot}}
\def\Hilbfunctor{\mathrm{Hilb}}
\def\Curvesstack{\mathcal{C}\!\mathit{urves}}
\def\Polarizedstack{\mathcal{P}\!\mathit{olarized}}
\def\Complexesstack{\mathcal{C}\!\mathit{omplexes}}
% \Pic is the operator that assigns to X its picard group, usage \Pic(X)
% \Picardstack_{X/B} denotes the Picard stack of X over B
% \Picardfunctor_{X/B} denotes the Picard functor of X over B
\def\Pic{\mathop{\mathrm{Pic}}\nolimits}
\def\Picardstack{\mathcal{P}\!\mathit{ic}}
\def\Picardfunctor{\mathrm{Pic}}
\def\Deformationcategory{\mathcal{D}\!\mathit{ef}}
\usepackage{cleveref}
\crefname{originallemma}{original lemma}{original lemmas}
\Crefname{originallemma}{Original lemma}{Original lemmas}
\makeatletter
\newlabel{schemes-lemma-characterize-closed-subspace}{{01HP}{source}{}{sourcepointer-1}{}}
\newlabel{schemes-lemma-characterize-closed-subspace@cref}{{[originallemma][0][]01HP}{[1][0][]source}{}{}{}}
\makeatother
\title{The universal property of the first order infinitesimal neighbourhood of an immersion}
\author{The Stacks Project Authors}
\date{stacks-project:more-morphisms.tex:local-snapshot-20261011; GNU FDL 1.2 or later}
\begin{document}
\maketitle
\noindent\textit{Editorial note.} All delivered mathematical text is the author's own, copied from the pinned Stacks Project chapter file more-morphisms.tex; no statement or proof step was rewritten, reordered or removed. The author writes 'We will only prove (1)'; that restriction is part of the author's text and is reproduced as written.\par\medskip
\noindent\textit{Editorial note.} The preamble is the chapter's own retained minimal preamble (sources/stacks-more-morphisms/retained-preamble.tex), a verbatim excerpt of the native Xy-pic setup and of the native theorem declarations and mathematical macros; the Xy-pic xy environment is required because the counted statement displays two commutative squares. No macro was invented or redefined.\par\medskip
\noindent\textit{Editorial note.} The cross-chapter standard prerequisite is located in the project's own tag list staged in the snapshot (sources/stacks-chow/tags.txt), at the verbatim line carrying its canonical tag.\par\medskip
\noindent\textit{Allowed standard prerequisites (editorial).}\begin{itemize}
\item The n-th order infinitesimal neighbourhood Z\_n of an immersion Z -> X. Same chapter, Definition environment labelled definition-first-order-infinitesimal-neighbourhood (source line 684, Stacks tag 04EY); the proof uses the open U of that construction and the ideal I cutting out Z in U. Supplies the objects Z' and Z\_n together with the ideal I and the fact that Z is closed in U; the factorisation argument is the counted content.
\item Criterion for a morphism to factor through the closed subscheme cut out by an ideal. Schemes, Lemma labelled schemes-lemma-characterize-closed-subspace (canonical Stacks tag 01HP); invoked at proof lines 739-741. Supplies the criterion that a morphism factors through the closed subscheme cut out by I as soon as the pullback of I becomes zero; the verification of that vanishing is retained in the counted span.
\end{itemize}\medskip
\ifdefined\hypertarget\hypertarget{sourcepointer-1}{}\fi
\noindent\textit{Source reference.} Original lemma 01HP: sources/stacks-chow/tags.txt, source lines 1950--1950. Cross-chapter standard prerequisite invoked once at proof lines 739-741 for the factorisation criterion through a closed subscheme cut out by an ideal; the reduction to the vanishing b\textasciicircum{}\#(b\textasciicircum{}\{-1\}(I\textasciicircum{}2)) = 0 is retained verbatim in the counted span.\par\medskip
\begin{lemma}
\label{lemma-first-order-infinitesimal-neighbourhood}
Let $i : Z \to X$ be an immersion of schemes.
\begin{enumerate}
\item The first order infinitesimal neighbourhood $Z'$ of $Z$ in $X$
has the following universal property: Given any commutative diagram
$$
\xymatrix{
Z \ar[d]_i & T \ar[l]^a \ar[d] \\
X & T' \ar[l]_b
}
$$
where $T \subset T'$ is a first order thickening over $X$, there exists
a unique morphism $(a', a) : (T \subset T') \to (Z \subset Z')$ of
thickenings over $X$.
\item For $n \geq 1$ the $n$th order infinitesimal neighbourhood $Z_n$
of $Z$ in $X$ has the following universal property: Given any commutative
diagram
$$
\xymatrix{
Z \ar[d]_i & T \ar[l]^a \ar[d] \\
X & T' \ar[l]_b
}
$$
where $T \subset T'$ is an $n$th order thickening over $X$, there exists
a unique morphism $(a', a) : (T \subset T') \to (Z \subset Z_n)$ of
thickenings over $X$.
\end{enumerate}
\end{lemma}

\begin{proof}
We will only prove (1).
Let $U \subset X$ be the open used in the construction of $Z'$, i.e., an
open such that $Z$ is identified with a closed subscheme of $U$ cut out by
the quasi-coherent sheaf of ideals $\mathcal{I}$.
Since $|T| = |T'|$ we see that $b(T') \subset U$. Hence we can
think of $b$ as a morphism into $U$. Let $\mathcal{J} \subset \mathcal{O}_{T'}$
be the ideal cutting out $T$. Since $b(T) \subset Z$ by the diagram above
we see that $b^\sharp(b^{-1}\mathcal{I}) \subset \mathcal{J}$. As
$T'$ is a first order thickening of $T$ we see that $\mathcal{J}^2 = 0$
hence $b^\sharp(b^{-1}(\mathcal{I}^2)) = 0$. By
Schemes, Lemma \ref{schemes-lemma-characterize-closed-subspace}
this implies that $b$ factors through $Z'$. Denote $a' : T' \to Z'$
this factorization and everything is clear.
\end{proof}

\end{document}
